\section{Software artifacts and reproducibility}
\label{appendix:software-and-reproducibility}

The software supporting this study comprises three repositories under the \texttt{kartAI} organization, run as a pipeline. \texttt{doppa-data-contribution} (\url{https://github.com/kartAI/doppa-data-contribution}) downloads the public source datasets, normalizes them to a shared GeoParquet schema, and publishes them to Blob Storage. \texttt{doppa} (\url{https://github.com/kartAI/doppa}) reads those datasets and runs the benchmarks, writing the per-iteration results back to Blob Storage. From those results, \texttt{doppa-analytics} (\url{https://github.com/kartAI/doppa-analytics}) generates the figures and tables reported in this thesis. Each repository documents its setup and run commands in its README.

Each repository is pinned to a specific commit: \texttt{doppa-data-contribution} at \texttt{6bea50e}, \texttt{doppa} at \texttt{a06e9fd}, and \texttt{doppa-analytics} at \texttt{30f1e39}. At these commits, \texttt{doppa-analytics} regenerates the reported figures and tables from the result data that \texttt{doppa} writes, and the output is copied into the thesis source. The repositories' READMEs give the full steps to reproduce the results from these commits.